\chapter{Reproduction Instructions}

\section{Environment Setup}

The SpriteCraft codebase requires Python 3.13 or later and PyTorch 2.0 or later. Dependencies are managed with \texttt{uv}:

\begin{lstlisting}[language=bash]
git clone https://github.com/RayZh-hs/SpriteCraft
cd SpriteCraft
uv sync
source .venv/bin/activate
\end{lstlisting}

\section{Preprocessing}

Place the vanilla Minecraft jar file and target resource pack archives (unmodified \texttt{.zip} or \texttt{.jar}) in \texttt{data/raw\_packs/}. Optionally create a manifest at \texttt{data/manifest.json} to specify pack roles:

\begin{lstlisting}[caption={Example manifest.json.}]
{
    "packs": [
        {
            "id": "vanilla",
            "archive": "server-1.21.jar",
            "role": "base",
            "style": "vanilla"
        },
        {
            "id": "Ashen_16x",
            "archive": "Ashen_16x.zip",
            "role": "train",
            "style": "stylized"
        }
    ]
}
\end{lstlisting}

Run preprocessing:

\begin{lstlisting}[language=bash]
spritecraft preprocess
\end{lstlisting}

This extracts archives, filters full-cube block textures by resolving blockstate models, resizes images to $32\times32$, and writes per-pack datasets to \texttt{data/processed/packs/<pack\_id>/}. A summary report is written to \texttt{data/processed/pack\_report.json}.

\section{Training}

Train the RecolorNet first (recommended, 2,500 steps per pack):

\begin{lstlisting}[language=bash]
spritecraft train --checkpoint-dir checkpoints/run33 --steps 2500 --mode recolor_only
\end{lstlisting}

Then train the diffusion model (10,000 steps per pack):

\begin{lstlisting}[language=bash]
spritecraft train --checkpoint-dir checkpoints/run33 --steps 10000 --recolor-checkpoint-dir checkpoints/run33
\end{lstlisting}

The \texttt{recolor\_first} mode trains RecolorNet then StyleAwareUNet sequentially for all packs. Monitor training with TensorBoard:

\begin{lstlisting}[language=bash]
tensorboard --logdir checkpoints/run33/tensorboard
\end{lstlisting}

\section{Generation}

Generate textures for a target pack:

\begin{lstlisting}[language=bash]
spritecraft generate --pack Ashen_16x --textures stone dirt oak_planks --run run33
\end{lstlisting}

Or sample randomly:

\begin{lstlisting}[language=bash]
spritecraft generate --pack Ashen_16x --random 8 --run run33
\end{lstlisting}

Outputs are written to \texttt{output/<pack\_id>/<bundle>/} containing the generated texture (\texttt{produced\_tex.png}), vanilla original (\texttt{original\_tex.png}), style references (\texttt{support.png}), comparison image (\texttt{comparison.png}), metrics (\texttt{metrics.json}), and metadata (\texttt{metadata.json}). If the target pack contains the queried texture, a source-of-truth image and MAE/pixel-accuracy metrics are included.

\section{Repository Layout}

\begin{lstlisting}[caption={Key project paths.}]
src/spritecraft/           Python package source
src/spritecraft/models/    Model definitions (unet, recolor, diffusion)
src/spritecraft/training/  Training loops and loss functions
src/spritecraft/inference/ Sampling and generation pipeline
src/spritecraft/data/      Preprocessing, dataset, support ranking
src/spritecraft/debug/     Runtime status and live inspection
data/raw_packs/            Input resource pack archives
data/processed/            Preprocessed per-pack datasets
checkpoints/<run>/         Training checkpoints and validation previews
output/<pack_id>/          Generated texture bundles
report/                    LaTeX report source
\end{lstlisting}

\section{Build Notes}

The report is configured to build with:

\begin{lstlisting}[language=bash]
make -C report
\end{lstlisting}

Generated files are written to \texttt{report/build/} and ignored by Git.

\section{Key Configuration Constants}

\begin{table}[htbp]
    \centering
    \caption{Key configuration constants in \texttt{src/spritecraft/config.py}.}
    \begin{tabular}{@{}ll@{}}
        \toprule
        Constant & Value \\
        \midrule
        IMAGE\_SIZE & 32 \\
        NUM\_TIMESTEPS & 40 \\
        MIN\_SUPPORT\_EXEMPLARS & 3 \\
        MAX\_SUPPORT\_EXEMPLARS & 3 \\
        VALIDATION\_FILENAMES & 20 fixed filenames \\
        MIN\_SHARED\_PACKS & 5 \\
        \bottomrule
    \end{tabular}
\end{table}
